\documentclass[10pt]{article}
\usepackage[margin=0.9in]{geometry}
\usepackage[T1]{fontenc}
\usepackage{lmodern,microtype,amssymb,enumitem,listings,xcolor}
\setlist{nosep,leftmargin=1.2em}
\lstset{basicstyle=\ttfamily\footnotesize,columns=fullflexible,keepspaces=true,frame=single,xleftmargin=0pt,
  aboveskip=4pt,belowskip=2pt,escapechar=|}
\newcommand{\K}[1]{K_{#1}}
\title{\vspace{-1.5em}Reasoning About Knowledge in TLA\textsuperscript{+}:\\Epistemic Properties for PlusCal Specs}
\author{A.\ Jesse Jiryu Davis\\MongoDB Research, New York}
\date{}
\begin{document}
\maketitle
\vspace{-2em}

\paragraph{Motivation.}
Engineers explain distributed protocols in terms of knowledge: a Raft leader should acknowledge a client's command only once it \emph{knows} that some follower has learned of the command; a client does not \emph{know} whether it is safe to retry a failed request; an eavesdropper should not \emph{learn} whether a login succeeded.
When it comes time to formally specify distributed protocols, however, we usually turn to temporal logic, e.g. the PlusCal dialect of TLA\textsuperscript{+}. Temporal logic says only how agents behave, not what they know. A marriage of PlusCal and epistemic logic~\cite{halpern-moses,fagin} would let us specify protocols more naturally, using the same concepts we use when describing them informally.

\paragraph{Approach.}
We propose a prototype that adds epistemic operators to PlusCal specs and checks epistemic properties.
Each PlusCal process is an agent, and its process-local variables are its local state.
TLC (the standard TLA\textsuperscript{+} model-checker) explores the spec and dumps its state graph to disk. Our new analysis tool then computes which states are indistinguishable to each agent.
It evalues epistemic formulae; for example, if a Raft leader is $L$, the set of followers is $F$, and the client's command is $\varphi$, then  ``the leader knows some follower knows the client's request'' is written $K_{L} S_{F} \varphi$.
Formulae live in spec comments, so existing tools are unaffected.
Temporal operators ($\Box$, $\Diamond$, $\leadsto$) can be combined with epistemic ones, and the quantifiers $\exists$ and $\forall$ are supported over finite literal domains.

\paragraph{Three new language features.}
A \emph{query} shows where a formula holds, and highlights those states in a visualization of the state graph.
A \emph{property} asserts a temporal-epistemic formula over the state graph, and reports violating states.
A \emph{precondition} checks that a labeled PlusCal step executes only at states in which the agent knows some fact.
These are demonstrated in the following annotations, from a Raft-like protocol:
\vspace{1em}
\begin{lstlisting}
\* At each state, query whether the leader L knows a follower i in the set F knows i has received a command:
\* KNOWLEDGE_QUERY     K(L, \E i \in F: K(i, received[i]))

\* Assert that sending a command to follower 1 leads to follower 1 receiving it:
\* KNOWLEDGE_PROPERTY  sent[1] ~> K(1, received[1])

\* Assert that a command is acknowledged only when the leader knows a follower i knows i received the command:
\* KNOWLEDGE_PRECONDITION AcknowledgeCommand:
\*     K(L, \E i \in F: K(i, received[i]))
\end{lstlisting}

\paragraph{Knowledge cannot be expressed in ordinary TLA\textsuperscript{+}.}
TLA\textsuperscript{+} expresses properties of individual \emph{behaviors}, which are linear paths through the state graph. But knowledge is a \emph{hyperproperty} involving states across many behaviors. An agent knows a fact $\varphi$ iff $\varphi$ is true in \emph{all} states that are indistinguishable from the current state, from the agent's perspective. These indistinguishable states may span many behaviors. Thus, no amount of shenanigans lets us express epistemic formulae in TLA\textsuperscript{+} or check them with TLC; we need an extension to the language itself.

We have noticed two useful categories of temporal-epistemic properties:
\vspace{1em}
\begin{itemize}
\item \textbf{Ignorance.} The requirement that an agent does \emph{not} know a fact. For example, if an agent $C$ eavesdrops on a login protocol, we might want to ensure that $C$ does not ever know whether the login succeeded or failed: $\Box(\lnot\K{C}\,\mathit{succeeded}\land\lnot\K{C}\lnot\mathit{succeeded})$.
\item \textbf{Necessity.} Whether a message is needed for an agent to learn something, e.g.\ that the leader does not know a follower has replicated a command, unless the follower sends an ack. This is useful for protocol minimization.
\end{itemize}

\paragraph{Limitations and related work.}
\begin{itemize}
\item \textbf{Related work:} MCMAS~\cite{mcmas} is an existing tool which checks computation tree logic extended with knowledge. MCMAS requires a dedicated modeling language and checker, whereas our prototype reuses PlusCal and TLC, which are the standard tools for distributed protocols.
\item \textbf{Scale:} Our evaluation is a post-processing pass over the full state graph, so it scales only as far as TLC can enumerate and store states; the examples are small.
\item \textbf{Theory:} We have not engaged deeply with the complexity results~\cite{bozzelli} or the symbolic techniques~\cite{lomuscio-penczek} for temporal-epistemic model checking.
\item \textbf{History:} Agents currently see only their local variables, so history must be added to the spec by hand.
\item \textbf{Quantifiers:} First-order quantification beyond finite literal domains is future research; evaluating a statement like ``$L$ knows there exists a follower $i$ in $F$ such that...' is complex if the set $F$ varies among states~\cite{fagin}.
\end{itemize}

\paragraph{Talk plan (30 minutes).}
\begin{enumerate}
\item We usually think and talk about protocols in terms of knowledge, but then we only specify behaviors. Example with Raft.
\item Epistemic logic in five minutes, including indistinguishability graphs.
\item Live demo: the three kinds of checks, using Raft and a login protocol.
\item Knowledge is a hyperproperty, so it cannot be expressed in TLA\textsuperscript{+}.
\item Questions for the audience: Which kinds of temporal-epistemic properties are most useful? What kind of protocols would benefit the most? How do we scale property-checking to specs with large state spaces?
\end{enumerate}

\begin{thebibliography}{9}\footnotesize
\bibitem{halpern-moses} J.~Halpern and Y.~Moses. Knowledge and common knowledge in a distributed environment. \emph{JACM} 37(3), 1990.
\bibitem{fagin} R.~Fagin, J.~Halpern, Y.~Moses, M.~Vardi. \emph{Reasoning About Knowledge}. MIT Press, 1995.
\bibitem{mcmas} A.~Lomuscio, H.~Qu, F.~Raimondi. MCMAS: a model checker for the verification of multi-agent systems. CAV 2009.
\bibitem{bozzelli} L.~Bozzelli, B.~Maubert, A.~Murano. On the complexity of model checking knowledge and time. \emph{ACM TOCL} 25(1), 2024.
\bibitem{lomuscio-penczek} A.~Lomuscio and W.~Penczek. Model checking temporal epistemic logic. In \emph{Handbook of Epistemic Logic}, ch.~8, College Publications, 2015.
\end{thebibliography}
\end{document}
